% ===== Document Preamble =====
\documentclass[11pt]{article}
\usepackage[margin=0.5in]{geometry}
\usepackage{amsmath,amssymb}
\usepackage{xcolor}
\usepackage{enumitem}
\usepackage{needspace}
\usepackage{comment}

% ===== Colors =====
\definecolor{AnsRed}{RGB}{170,0,0}
\definecolor{WrongBlue}{RGB}{0,0,160}
\definecolor{SkipGray}{RGB}{120,120,120}

% ===== Choice list styling (tight, consistent) =====
\setlist[enumerate]{%
  leftmargin=1.5em,
  itemsep=0.25em,
  topsep=0.25em,
  label=\textbf{\Alph*.}
}

% ===== Style Wrappers =====
\newcommand{\QuestionStyle}[1]{{\large\bfseries\boldmath #1}}
\newcommand{\AnswerStyle}[1]{{\large\bfseries\boldmath #1}}
\newcommand{\ReasonStyle}[1]{{\normalsize #1}}

% ===== Show / Hide Reasoning Toggle =====
\newif\ifshowreason
\showreasontrue      % Instructor version
% \showreasonfalse   % Student version

% ===== Choice Item Macros (Answer + Reasoning) =====
% Usage:
%   \Correct{Answer}{Reasoning}
%   \Wrong{Answer}{Reasoning}

\newcommand{\Correct}[2]{%
  \item \AnswerStyle{\textcolor{AnsRed}{#1}}%
  \ifshowreason\\[-0.2em]
    \ReasonStyle{\textcolor{AnsRed}{#2}}%
  \fi
}

\newcommand{\Wrong}[2]{%
  \item \AnswerStyle{\textcolor{WrongBlue}{#1}}%
  \ifshowreason\\[-0.2em]
    \ReasonStyle{\textcolor{WrongBlue}{#2}}%
  \fi
}

% ===== Multiple-Choice Question Block Macro =====
% Usage: \QBlock{Q1}{stem...}{<items>}{answer key}

\newcommand{\QBlock}[4]{%
  \Needspace{8\baselineskip}%
  \noindent
  \begin{minipage}{\linewidth}
    \QuestionStyle{#1.~#2}

    \vspace{0.25em}
    \begin{enumerate}
      #3
    \end{enumerate}

    \vspace{0.1em}
    \noindent\textcolor{AnsRed}{\textit{\textbf{Answer:}~#4}}

    \vspace{0.25em}
    \hrule height 0.4pt
  \end{minipage}
  \vspace{0.75em}
}

% ===== Open-Ended Question Block Macro =====
% Usage: \OEBlock{Q1}{stem...}{answer text}

\newcommand{\OEBlock}[3]{%
  \Needspace{5\baselineskip}%
  \noindent
  \begin{minipage}{\linewidth}
    \QuestionStyle{#1.~#2}

    \vspace{0.25em}
    \noindent\textcolor{AnsRed}{\textit{\textbf{Answer:}~#3}}

    \vspace{0.25em}
    \hrule height 0.4pt
  \end{minipage}
  \vspace{0.75em}
}

% ===== Multi-Part Question Block Macro =====
% Usage: \MPBlock{Q1}{stem...}{<parts>}

\newcommand{\MPBlock}[3]{%
  \Needspace{8\baselineskip}%
  \noindent
  \begin{minipage}{\linewidth}
    \QuestionStyle{#1.~#2}

    \vspace{0.25em}
    #3

    \vspace{0.25em}
    \hrule height 0.4pt
  \end{minipage}
  \vspace{0.75em}
}

% ===== MC Sub-Part Macro =====
% Usage: \MCPart{a}{stem...}{<items>}{answer key}

\newcommand{\MCPart}[4]{%
  \par\vspace{0.35em}
  \noindent\textbf{(#1)}~#2
  \begin{enumerate}
    #3
  \end{enumerate}
  \vspace{0.1em}
  \noindent\textit{\textbf{Answer:}~#4}
  \par
}

% ===== Open-Ended Sub-Part Macro =====
% Usage: \OEPart{a}{stem...}{answer text}

\newcommand{\OEPart}[3]{%
  \par\vspace{0.35em}
  \noindent\textbf{(#1)}~#2
  \par\vspace{0.15em}
  \noindent\textit{\textbf{Answer:}~#3}
  \par
}

% ===== Skipped Question Note Macro =====
% Usage: \SkipBlock{Q1}{reason}

\newcommand{\SkipBlock}[2]{%
  \noindent\textcolor{SkipGray}{\textit{[#1 omitted: #2]}}%
  \vspace{0.5em}\par
}

% ===== Document Start =====

\begin{document}

\begin{center}
{\huge \textbf{Physics Isomorphic Questions}}\\[0.25em]
{\Large Vectors and One-Dimensional Kinematics}\\[0.25em]
{\large Questions generated below. Multiple-choice answers in RED; incorrect answers in BLUE.}
\end{center}

\vspace{0.75em}

\QBlock{Q1}{A survey drone starts at the origin of a coordinate grid and detects a beacon at position \(\vec r = -7\hat i + 24\hat j\), where distances are measured in kilometers. Which of the following statements is correct?}{
\Wrong{The distance from the drone to the beacon is \(31\) km.}{This comes from adding the magnitudes of the components, \(7+24\), instead of finding the vector magnitude. Distance from the origin must be computed with \(|\vec r|=\sqrt{x^2+y^2}\), not by simple addition.}
\Correct{The distance from the drone to the beacon is \(25\) km.}{The beacon's position vector has components \(-7\) km and \(24\) km, so its distance from the origin is the magnitude \(|\vec r|=\sqrt{(-7)^2+24^2}=\sqrt{49+576}=\sqrt{625}=25\,\mathrm{km}\). The sign of the \(x\)-component affects direction but not the distance, so the correct choice is \(25\) km.}
\Wrong{The distance from the drone to the beacon is \(24\) km.}{This treats only the vertical component as the full distance. Both the \(x\)- and \(y\)-components contribute to the magnitude of the position vector.}
\Wrong{The distance from the drone to the beacon is \(7\) km.}{This uses only the horizontal component and ignores the vertical displacement. A vector's magnitude must include all perpendicular components through the Pythagorean relation.}
}{B}

\QBlock{Q2}{The position of a particle moving along the \(x\)-axis is given by \(\vec s(t)=\big((0.80\,\mathrm{m/s^3})t^3-(1.90\,\mathrm{m/s^2})t^2\big)\hat i\). Which of the following statements is correct?}{
\Correct{The magnitude of the acceleration at \(t=0\) is \(3.8\,\mathrm{m/s^2}\).}{Differentiate twice: \(\vec v(t)=\big(2.40t^2-3.80t\big)\hat i\) and \(\vec a(t)=\big(4.80t-3.80\big)\hat i\). At \(t=0\), the acceleration is \(-3.80\hat i\,\mathrm{m/s^2}\), so its magnitude is \(3.8\,\mathrm{m/s^2}\).}
\Wrong{This motion has constant acceleration.}{This error comes from overlooking the cubic term in position. Since \(\vec a(t)=\big(4.80t-3.80\big)\hat i\), the acceleration depends on time and is not constant.}
\Wrong{The magnitude of the acceleration is \(1.9\,\mathrm{m/s^2}\).}{This confuses the coefficient of the \(t^2\) term in displacement with acceleration itself. Acceleration must be obtained from the second derivative of the displacement function.}
\Wrong{The magnitude of the acceleration at \(t=0\) is \(0\,\mathrm{m/s^2}\).}{This mistake assumes the acceleration must vanish initially. Evaluating \(\vec a(t)=\big(4.80t-3.80\big)\hat i\) at \(t=0\) gives a nonzero value.}
}{A}

\QBlock{Q3}{A robot moves from position \(\vec r_A = 1\hat i - 2\hat j\) m to position \(\vec r_B = 7\hat i + 6\hat j\) m in \(5.0\) s. Which of the following statements is correct?}{
\Wrong{The robot must have traveled a distance of \(10\) m.}{The value \(10\,\mathrm{m}\) is the magnitude of the displacement, not necessarily the path length traveled. The actual distance traveled depends on the path and cannot be determined from only the initial and final positions.}
\Wrong{The robot's speed at \(t=0\) must have been \(2.0\,\mathrm{m/s}\).}{This comes from dividing displacement magnitude by elapsed time and interpreting that as the initial speed. Endpoint positions over a time interval do not uniquely determine the speed at the starting instant.}
\Correct{The magnitude of the displacement is \(10\) m.}{The displacement is \(\Delta\vec r=\vec r_B-\vec r_A=(7-1)\hat i + \big(6-(-2)\big)\hat j = 6\hat i+8\hat j\,\mathrm{m}\). Its magnitude is \(|\Delta\vec r|=\sqrt{6^2+8^2}=\sqrt{100}=10\,\mathrm{m}\), so this is the only statement that must be true.}
\Wrong{The robot's acceleration at \(t=0\) must have been \(0.40\,\mathrm{m/s^2}\).}{This incorrectly assumes acceleration can be inferred from displacement and total time alone. Without more information about the motion or velocity, the initial acceleration is not determined.}
}{C}

\end{document}
